\begin{papercard}{cohen1960}
\cf{Stages} \stagetag{I}. The original paper proposing $\kappa$, a chance-corrected agreement coefficient for two judges classifying units into unordered categories. It states the assumptions, derives the coefficient, its limits and maximum, its standard error and significance test. \emph{Note:} the PDF is a scan with OCR errors, so snippets reproduce the printed text layer.

\cf{Problem and setting} Nominal classification is a complex judgement by a ``two-legged meter'', so one must determine how reproducible the judgements are, by having two or more judges classify a sample independently \pc{cohen1960}{1}{it becomes important to determine the extent to which these judg}. The analogy is to the coefficient of equivalence: judges are like alternate forms.

\cf{Assumptions} (1) The units are independent. (2) The categories are independent, mutually exclusive and exhaustive. (3) The judges operate independently. There is no criterion of correctness, the judges are deemed equally competent, and the distribution of judgments over categories is unrestricted. All discrepancies count equally (no ordering) \pc{cohen1960}{2}{The units are independent}.

\cf{Why $\chi^2$ and percent agreement fail} The crude approach, the proportion of agreement $p_o$, is inadequate because some agreement is expected by chance (the product of the marginals). $\chi^2$ or the contingency coefficient $C$ tests \emph{association}, not agreement, and is inflated by any departure from chance, including disagreement. In the example, $p_o=.29$ is below chance agreement $p_c=.35$, yet $\chi^2=64.59$ is highly significant: ``the judgments are associated, but unfortunately not in the direction of agreement'' \pc{cohen1960}{3}{It is readily demonstrable that the use ofx}.

\cf{The coefficient} $\kappa=(p_o-p_c)/(1-p_c)$, with $p_c=\sum p_{iA}p_{iB}$. It is ``the proportion of chance-expected disagreements which do not occur'', or the proportion of agreement after chance agreement is removed. Related coefficients: Guttman's $\lambda$ (Goodman \& Kruskal) uses modal categories, which is questionable for reliability since each judge distributes judgments over all $k$ categories. Scott's $\pi$ assumes a known population distribution identical for both judges. Cohen objects that a proclivity to distribute judgments differently over categories is itself a source of disagreement, and so keeps each judge's own marginals \pc{cohen1960}{4}{The determination of pc for K follows from the logic of the reliability situation}, \pc{cohen1960}{5}{since one source of disagreement between a pair of judges is precisely their proclivity to distribute their judgments differently}.

\cf{Limits and maximum} $\kappa=0$ when agreement equals chance, $+1$ only with perfect agreement, which requires identical marginals. The lower limit depends on the marginals (between 0 and $-1$) \pc{cohen1960}{5}{For perfect agreement there is a necessary condition that the p}. The \emph{maximum} $\kappa_M$ allowed by the marginals is computed by pairing the smaller marginal proportions (example: $p_{oM}=.90$, $\kappa_M=.85$); $1-\kappa_M$ indicates ``the fuzziness of his category boundaries'' and can guide judge training. But disagreement forced by marginal disagreement ``is disagreement'' \pc{cohen1960}{6}{is the maximum value of K permitted by the marginals}. For two categories with equal marginals $\kappa=\phi$.

\cf{Sampling and testing} Approximate standard error, treating $p_c$ as constant (adequate for $N>100$), normal confidence limits $\kappa\pm1.96\sigma_\kappa$, a test for the difference between two independent $\kappa$s, and a test of $\kappa=0$. Cohen adds that testing against zero is ``of as little value'' as for any reliability coefficient, since a much better reliability is expected \pc{cohen1960}{8}{to know merely that K is beyond chance is trivial since one usually expects much more than}. Illustration: $\kappa=.492$, $\kappa_M=.831$, 95\,\% limits .384--.600 \pc{cohen1960}{8}{With K = .492, we see that just under half of the joint judgments are agreements}.

\cf{Use for this thesis} The primary source for the definition and assumptions of $\kappa$ (independent units, mutually exclusive and exhaustive categories, independent judges), and for the calculation of $\kappa_{\max}$ as a diagnostic of category-boundary fuzziness. Because Cohen deliberately keeps each judge's marginals in the chance term, later critics (Krippendorff, Gwet, Feinstein \& Cicchetti, Hallgren) disagree with this choice. The thesis should therefore cite Cohen for the formula and assumptions, and a comparative source (Hallgren; Krippendorff) for choosing among $\kappa$, $\pi$ and $\alpha$. Cohen does not give interpretation bands or sample-size rules.
\end{papercard}
